% GENERATED FILE. Edit canonical lessons, then run npm run generate:books.
% canonical-source-hash: fnv1a64:df939bf0a5fa7bbd
% canonical-lessons: TA-C245-cumai, TA-C245-kural, TA-C245-uraiyatal, TA-C245-an, TA-C245-pen

\chapter{Load, Voice, Conversation, Man, Woman}
\label{ch:ta-load-voice-conversation-man-woman}
\hlchaptermodality{245}

\begin{chapteropening}
\textbf{By the end of this chapter:} \emph{I can say a load, a voice, a conversation, a man and a woman.}

Complete the last lesson of chapter 245: I can say a load, a voice, a conversation, a man and a woman.
\end{chapteropening}

\section[\texorpdfstring{\ta{சுமை} (cumai)}{சுமை (cumai)}]{\ta{சுமை} (cumai) — a load, a burden}

\label{lesson:TA-C245-cumai}



\begin{quote}
Before the new one: say the Tamil for a document, then the Tamil for a passport.
\end{quote}

\subsection*{You'll want to know: \ta{சுமை}}
\textbf{\ta{சுமை}} — \emph{cumai} — \textquotedblleft{}a load, a burden\textquotedblright{}.

\begin{grammarlens}[title={What you've built}]
One more word for this chapter.
\end{grammarlens}

\subsection*{Your turn}
\begin{itemize}
\raggedright
  \item \emph{Say it:} \emph{cumai}
  \item \emph{Say it:} \emph{cumai}, once more
  \item \emph{Say it:} say \emph{cumai}
  \item \emph{Recall:} say the Tamil for a document, then the Tamil for a passport, then say \emph{cumai} again
\end{itemize}

\begin{tcolorbox}[breakable,colback=teal!4,colframe=teal!35!black,title={Before you move on}]
What does \ta{சுமை} mean? (\textquotedblleft{}A load, a burden\textquotedblright{}.) Say it once more.
\end{tcolorbox}
\section[\texorpdfstring{\ta{குரல்} (kural)}{குரல் (kural)}]{\ta{குரல்} (kural) — a voice}

\label{lesson:TA-C245-kural}



\begin{quote}
Before the new one: say the Tamil for a passport, then the Tamil for a load.
\end{quote}

\subsection*{You'll want to know: \ta{குரல்}}
\textbf{\ta{குரல்}} — \emph{kural} — \textquotedblleft{}a voice\textquotedblright{}.

\begin{grammarlens}[title={What you've built}]
One more word for this chapter.
\end{grammarlens}

\subsection*{Your turn}
\begin{itemize}
\raggedright
  \item \emph{Say it:} \emph{kural}
  \item \emph{Say it:} \emph{kural}, once more
  \item \emph{Say it:} say \emph{kural}
  \item \emph{Recall:} say the Tamil for a passport, then the Tamil for a load, then say \emph{kural} again
\end{itemize}

\begin{tcolorbox}[breakable,colback=teal!4,colframe=teal!35!black,title={Before you move on}]
What does \ta{குரல்} mean? (\textquotedblleft{}A voice\textquotedblright{}.) Say it once more.
\end{tcolorbox}
\section[\texorpdfstring{\ta{உரையாடல்} (uraiyāṭal)}{உரையாடல் (uraiyāṭal)}]{\ta{உரையாடல்} (uraiyāṭal) — a conversation}

\label{lesson:TA-C245-uraiyatal}



\begin{quote}
Before the new one: say the Tamil for a load, then the Tamil for a voice.
\end{quote}

\subsection*{You'll want to know: \ta{உரையாடல்}}
\textbf{\ta{உரையாடல்}} — \emph{uraiyāṭal} — \textquotedblleft{}a conversation\textquotedblright{}.

\begin{grammarlens}[title={What you've built}]
One more word for this chapter.
\end{grammarlens}

\subsection*{Your turn}
\begin{itemize}
\raggedright
  \item \emph{Say it:} \emph{uraiyāṭal}
  \item \emph{Say it:} \emph{uraiyāṭal}, once more
  \item \emph{Say it:} say \emph{uraiyāṭal}
  \item \emph{Recall:} say the Tamil for a load, then the Tamil for a voice, then say \emph{uraiyāṭal} again
\end{itemize}

\begin{tcolorbox}[breakable,colback=teal!4,colframe=teal!35!black,title={Before you move on}]
What does \ta{உரையாடல்} mean? (\textquotedblleft{}A conversation\textquotedblright{}.) Say it once more.
\end{tcolorbox}
\section[\texorpdfstring{\ta{ஆண்} (āṇ)}{ஆண் (āṇ)}]{\ta{ஆண்} (āṇ) — a man, male}

\label{lesson:TA-C245-an}



\begin{quote}
Before the new one: say the Tamil for a voice, then the Tamil for a conversation.
\end{quote}

\subsection*{You'll want to know: \ta{ஆண்}}
\textbf{\ta{ஆண்}} — \emph{āṇ} — \textquotedblleft{}a man, male\textquotedblright{}.

\begin{grammarlens}[title={What you've built}]
One more word for this chapter.
\end{grammarlens}

\subsection*{Your turn}
\begin{itemize}
\raggedright
  \item \emph{Say it:} \emph{āṇ}
  \item \emph{Say it:} \emph{āṇ}, once more
  \item \emph{Say it:} say \emph{āṇ}
  \item \emph{Recall:} say the Tamil for a voice, then the Tamil for a conversation, then say \emph{āṇ} again
\end{itemize}

\begin{tcolorbox}[breakable,colback=teal!4,colframe=teal!35!black,title={Before you move on}]
What does \ta{ஆண்} mean? (\textquotedblleft{}A man, male\textquotedblright{}.) Say it once more.
\end{tcolorbox}
\section[\texorpdfstring{\ta{பெண்} (peṇ)}{பெண் (peṇ)}]{\ta{பெண்} (peṇ) — a woman, female}

\label{lesson:TA-C245-pen}



\begin{quote}
Before the new one: say the Tamil for a conversation, then the Tamil for a man.
\end{quote}

\subsection*{You'll want to know: \ta{பெண்}}
\textbf{\ta{பெண்}} — \emph{peṇ} — \textquotedblleft{}a woman, female\textquotedblright{}.

\begin{grammarlens}[title={What you've built}]
One more word for this chapter.
\end{grammarlens}

\subsection*{Your turn}
\begin{itemize}
\raggedright
  \item \emph{Say it:} \emph{peṇ}
  \item \emph{Say it:} \emph{peṇ}, once more
  \item \emph{Say it:} say \emph{peṇ}
  \item \emph{Recall:} say the Tamil for a conversation, then the Tamil for a man, then say \emph{peṇ} again
\end{itemize}

\begin{tcolorbox}[breakable,colback=teal!4,colframe=teal!35!black,title={Before you move on}]
What does \ta{பெண்} mean? (\textquotedblleft{}A woman, female\textquotedblright{}.) Say it once more.
\end{tcolorbox}
